La glissade doit absolument passer par 4 points fixes ($A$,$B$,$C$,$D$) et un
point dont la hauteur est variable entre 10 et 15 mètres ($E$). Si
les coordonés exactes du point $E$ étaient connues à l'avance, créer un fonction
polynomiale passant par les 5 points pourrait être résolue en utilisant les
coordonés pour développer un système d'équations avec 5 équations et 4 inconnus
(polynomiale d'ordre 3). Plusieures suppositions sont émises dans Matlab,
chacunes avec une valeur de $E$ différentes et uniformément réparties dans la
plage permise. Pour chaque système d'équation, les coefficients de la
polynomiale sont résolues. Afin d'être bref seul la résolution pour le $E$
choisi est montrée puisque la démarche est toujours la même (le choix est justifié
plus tard):

\begin{align}
	F(x) &= a_0 + a_1 x^1 + a_2 x^2 + a_3 x^3 \\
\begin{bmatrix} 30\\ 20\\ 20\\ 16\\ 12.5 \end{bmatrix}
&=
\begin{bmatrix}
1 & 0 & 0^2 & 0^3 & 0^4\\
1 & 7 & 7^2 & 7^3 & 7^4\\
1 & 15 & 15^2 & 15^3 & 15^4\\
1 & 20 & 20^2 & 20^3 & 20^4\\
1 & 25 & 25^2 & 25^3 & 25^4
\end{bmatrix}
\begin{bmatrix} a_0\\ a_1\\ a_2\\ a_3\\ a_4 \end{bmatrix} \Rightarrow
\begin{bmatrix} a_0\\ a_1\\ a_2\\ a_3\\ a_4 \end{bmatrix}
\end{align}


